\documentclass[11pt]{article}
\usepackage[T1]{fontenc}
\usepackage{lmodern,amsmath,amssymb,amsthm,booktabs}
\usepackage[a4paper,margin=27mm]{geometry}
\usepackage[hidelinks]{hyperref}
\usepackage{microtype}
\newtheorem{theorem}{Theorem}[section]
\newtheorem{lemma}[theorem]{Lemma}
\newtheorem{corollary}[theorem]{Corollary}
\newtheorem{proposition}[theorem]{Proposition}
\theoremstyle{remark}
\newtheorem{remark}[theorem]{Remark}
\newcommand{\Th}{\Theta}
\newcommand{\bn}[1]{\binom{#1}{#1/2}}
\newcommand{\N}{\mathcal N}
\title{Exact facet maximizers for simple theta graphs and the bicyclic extremal problem}
\author{mio-qwq/math project\\\small STRUCT2-20261012 research manuscript}
\date{13 October 2026}
\begin{document}
\maketitle
\begin{abstract}
We determine the unique unordered path-length triple maximizing the number
of facets of the symmetric edge polytope of a simple theta graph, for every
order at which this class is nonempty. The maximizing exceptional-parity
path has length one at even order and length two at odd order; its two
remaining paths are as balanced as their parity permits. This proves the
simple-graph version of Braun--Bruegge Conjecture 33 and strengthens its
stated range and conclusion. Exact ratios to their bicyclic candidate
bound are strictly less than one at every order at least five, proving
Conjecture 32 on its natural domain. A block reduction then gives the
maximum for every connected simple graph with one more edge than vertices,
together with the maximizing block data. The argument treats length-one
paths explicitly and uses a joint bound on both length-weighted counting
terms. Its universal conclusions rest on exact inequalities, with separate
definition-level computational checks supplied for reproducibility.
\end{abstract}

\section{Problem, domain, and results}
For a finite connected simple graph $G$ on $[n]$, let
\[
 P_G=\operatorname{conv}\{\pm(e_i-e_j):ij\in E(G)\},\qquad
 \N(G)=\#\{\text{facets of }P_G\}.
\]
A theta graph $\Th(a,b,c)$ consists of three internally vertex-disjoint
paths with common distinct endpoints, of edge lengths $a,b,c$. We work
throughout in the domain
\begin{equation}\label{eq:domain}
 a,b,c\ge1,\quad \#\{i:a_i=1\}\le1,\quad a+b+c=n+1.
\end{equation}
Thus it has $n$ vertices and $n+1$ edges, and $n\ge4$. Two purported
length-one paths would duplicate the same edge of a simple graph.
Braun and Bruegge use the notation $\mathrm{CB}(a,b,c)$ for this graph
\cite{BB}. Their journal Conjecture 33 proposes four maximizers for
$n\ge10$; Conjecture 32 asks for the remaining mixed-parity bound in
their bicyclic extremal problem.

\begin{theorem}\label{thm:theta}
For every $n\ge4$, the maximizing unordered triple in \eqref{eq:domain}
is unique and is given by the following table.
\begin{center}
\begin{tabular}{ccc}\toprule
Order & Parity of $k$ & Maximizing triple\\\midrule
$n=2k-1$ & even & $(k-1,k-1,2)$\\
$n=2k-1$ & odd & $(k,k-2,2)$\\
$n=2k$ & even & $(k,k,1)$\\
$n=2k$ & odd & $(k+1,k-1,1)$\\\bottomrule
\end{tabular}
\end{center}
The odd-order rows have $k\ge3$; the even-order rows have $k\ge2$.
\end{theorem}

For odd $s\ge3$, put $h(s)=s\binom{s-1}{(s-1)/2}$, and, for $k\ge3$,
\begin{equation}\label{eq:M}
 W(k)=\begin{cases}h(k-1)h(k+1),&k\text{ even},\\h(k)^2,&k\text{ odd},\end{cases}
 \qquad M(2k-1)=W(k),\quad M(2k)=2W(k).
\end{equation}
This is the graph-defined candidate of \cite[Definition 12]{BB}.

\begin{theorem}\label{thm:bicyclic}
For every $n\ge5$, the maximum of $\N(G)$ among connected simple graphs
with $n$ vertices and $n+1$ edges is $M(n)$. Every such graph with a theta
core has strictly fewer facets. A maximizing graph has two odd cycle
blocks with the balanced lengths in \eqref{eq:M}, and zero bridges at
odd order or exactly one bridge at even order.
\end{theorem}

At even order, the maximizing block data may give different graph
isomorphism types: the bridge can join disjoint cycles or be a pendant
edge on their one-vertex wedge. We do not assert unique whole-graph
maximizers there. At $n=4$, the sole eligible graph is $K_4-e$ and its
facet count is 12. At $n=3$, no simple graph has four edges. The expression
in \cite[Definition 12]{BB} contains an undefined $C_1$ at these small
orders, so we do not introduce a one-cycle convention.

\section{The counting formula and its boundary cases}
We use the facet-labeling characterization of Higashitani, Jochemko and
Micha\l ek \cite{HJM}, recalled in \cite[Theorem 5]{BB}. After one vertex
label is normalized to zero, facets correspond exactly to integer
functions $f$ whose difference has absolute value at most one on every
edge and whose edges of absolute difference one form a connected
spanning subgraph. Opposite functions are not identified.

For $m\ge0$, define
\[
 B_m(t)=\begin{cases}
 \binom m{(m+t)/2},& |t|\le m,\ m+t\equiv0\pmod2,\\
 0,&\text{otherwise},
 \end{cases}
 \qquad T(a,b,c)=\sum_tB_a(t)B_b(t)B_c(t).
\]
In particular $B_0(t)=1[t=0]$; this is a counting row and does not
introduce a zero-length graph path.

\begin{lemma}\label{lem:count}
For common-parity lengths, $\N(\Th(a,b,c))=T(a,b,c)$. For mixed parity,
name the same-parity lengths $x,y$ and the other length $z$. Then
\begin{equation}\label{eq:weighted}
 \N(x,y,z):=\N(\Th(x,y,z))
 =zT(x,y,z-1)+xyT(x-1,y-1,z).
\end{equation}
This identity includes every legal length-one boundary in
\eqref{eq:domain}.
\end{lemma}
\begin{proof}
Normalize one common endpoint to zero and let the other endpoint label
be $t$. Along each path increments belong to $\{-1,0,1\}$. Two zero
increments on one path disconnect the segment between them in the
saturated subgraph, so there is at most one zero on each path. At least
one path has no zero, since otherwise the two endpoints are disconnected
in that subgraph. Conversely, these two conditions make the saturated
subgraph connected and spanning.

For equal parities, equality of endpoint sums forces all three paths to
have the same zero parity; only zero zeros on every path is allowed.
This gives $T$. For mixed parity the two possibilities are: the $z$-path
has one zero and the $x,y$ paths have none; or the $z$-path has none and
each of the other two has one. Choosing zero positions gives the
coefficients $z$ and $xy$ in \eqref{eq:weighted}. With a length-one path,
deleting its zero increment leaves the row $B_0$. Thus the same argument
covers the boundary rather than substituting an undefined path into a
positive-length formula. These counts agree with the separate clauses
of \cite[Corollary 25]{BB}.
\end{proof}

For same-parity $p,q\ge0$, Vandermonde's identity gives
\begin{equation}\label{eq:vander}
 \sum_t B_p(t)B_q(t)=B_{p+q}(0),\qquad
 T(a,b,c)\le \max_tB_a(t)\,B_{b+c}(0).
\end{equation}
Write $b_j=\binom{2j}{j}$ for $j\ge0$. Its successive ratio is
\begin{equation}\label{eq:rho}
 \rho_j=b_{j+1}/b_j=4-2/(j+1),
\end{equation}
which is positive and strictly increasing. It follows that, for $u\ge0$
and $v\ge1$,
\begin{equation}\label{eq:product}
 \frac{b_ub_v}{b_{u+v}}\le\frac{u+1}{2u+1},\qquad
 b_rb_s\le2b_{r+s-1}\quad(r,s\ge1).
\end{equation}
Indeed the first quotient is nonincreasing in $v$ because its successive
ratio is $\rho_v/\rho_{u+v}\le1$; evaluate at $v=1$. The second quotient
equals 2 at $r=1$ and has successive ratio $\rho_r/\rho_{r+s-1}\le1$.

\begin{lemma}\label{lem:short}
For even $p,q\ge2$ and odd $x,y\ge1$ with at most one equal to one,
\begin{align}
 \N(p,q,1)&=\frac{pq+2}{2}B_p(0)B_q(0),\label{eq:short1}\\
 \N(x,y,2)&=\frac{xy+5}{4}B_{x+1}(0)B_{y+1}(0).\label{eq:short2}
\end{align}
\end{lemma}
\begin{proof}
For \eqref{eq:short1}, only $t=0$ occurs in the first term of
\eqref{eq:weighted} and only $t=\pm1$ in the second; use
$B_{p-1}(1)=B_p(0)/2$.
For \eqref{eq:short2}, put $A_x=B_{x-1}(0)$ and
$r_x=(x-1)/(x+1)$, and similarly for $y$. Then
\[
 B_x(1)=\frac{2x}{x+1}A_x,\quad B_{x-1}(2)=r_xA_x,\quad
 B_{x+1}(0)=\frac{4x}{x+1}A_x.
\]
These also hold at $x=1$. The two terms in \eqref{eq:weighted} give
\[
 \N(x,y,2)=4B_x(1)B_y(1)+2xyA_xA_y(1+r_xr_y)
 =\frac{4xy(xy+5)}{(x+1)(y+1)}A_xA_y.
\]
This is \eqref{eq:short2}. The length multipliers and the constant 5
are essential.
\end{proof}

\section{Strict compression of the exceptional-parity path}
\begin{proposition}\label{prop:even}
For even $x,y\ge2$ and odd $z\ge1$,
\[
 \N(x,y,z)\le\N(x,y+z-1,1),
\]
strictly when $z>1$.
\end{proposition}
\begin{proof}
Put $D=B_{y+z-1}(0)>0$. The peaks of the even and odd rows of lengths
$x,x-1$ are $B_x(0)$ and $B_x(0)/2$. Jointly applying
\eqref{eq:vander} to \eqref{eq:weighted} gives
\[
 \N(x,y,z)\le B_x(0)D\left(z+\frac{xy}{2}\right).
\]
By \eqref{eq:short1}, the compressed count is
$B_x(0)D(1+x(y+z-1)/2)$. The difference of these factors is
$(z-1)(x-2)/2\ge0$. It is positive if $z>1$ and $x>2$.
If $z>1$ and $x=2$, the first envelope itself is strict: at $t=2$ its
weight $B_y(2)B_{z-1}(2)$ is positive, whereas $B_2(2)=1<B_2(0)=2$.
This covers the boundary $x=y=2$. If $z=1$, the graph is unchanged.
\end{proof}

\begin{proposition}\label{prop:odd}
For odd $x,y\ge1$, even $z\ge2$, and at most one of $x,y$ equal to one,
\[
 \N(x,y,z)\le\N(x,y+z-2,2),
\]
after naming $x\ge3$. The inequality is strict when $z>2$.
\end{proposition}
\begin{proof}
Set
\[
 A=B_{x-1}(0),\quad D=B_{y+z-1}(0),\quad
 E=B_{y-1}(0)B_z(0),\quad r=(x-1)/(x+1).
\]
For every even $t$, the largest noncentral entry of the row $x-1$ is
at $\pm2$, so
\[
 B_{x-1}(t)\le rA+(1-r)A1[t=0].
\]
Vandermonde now bounds the second counting term by
$T(x-1,y-1,z)\le rAD+(1-r)AE$. The first term has odd-row peak
$2xA/(x+1)$. Together they give
\begin{equation}\label{eq:joint}
 \N(x,y,z)\le\frac{xA}{x+1}
 \bigl((2z+y(x-1))D+2yE\bigr).
\end{equation}
Take $u=(y-1)/2$ and $v=z/2$ in \eqref{eq:product}. This yields
$E/D\le(y+1)/(2y)$, including $y=1$. Equation \eqref{eq:joint}
therefore reduces to
\[
 \N(x,y,z)\le\frac{xAD}{x+1}(xy+2z+1).
\]
Write $y'=y+z-2$, positive and odd. By \eqref{eq:short2}, the compressed
count is $xAD(xy'+5)/(x+1)$. The factor difference is
\[
 xy'+5-(xy+2z+1)=(x-2)(z-2).
\]
It is positive for $z>2$ and vanishes at the unchanged case $z=2$.
Both compressed triples preserve total length and the simple domain.
\end{proof}

\section{Strict balancing and exclusion of common parity}
Let $u,v\ge2$ be even and $u\ge v+4$. Put
\[
 P=uv,\quad S=u+v,\quad d=u-v-2>0,\quad
 U=u(v+1),\quad V=(u-1)(v+2)=U+d.
\]
The inward transfer $(u,v)\mapsto(u-2,v+2)$ increases the product to
$P+2d$. Its central-binomial product ratio is
\[
 \frac{B_{u-2}(0)B_{v+2}(0)}{B_u(0)B_v(0)}=\frac UV.
\]
For the length-one case \eqref{eq:short1}, cross multiplication gives
\[
 U(P+2d+2)-V(P+2)=d(uv+2u-2)>0.
\]
For the length-two case, write $x=u-1,y=v-1$. The nonconstant weight
in \eqref{eq:short2} is $P-S+6$, and its comparison gives
\[
 U(P+2d-S+6)-V(P-S+6)=d(uv+3u+v-6)>0.
\]
Thus the remaining pair is uniquely as balanced as parity permits.
The transfer keeps positive lengths and at most one length one.

It remains to exclude triples of common parity. For odd vertex count,
all three lengths are even. Write $a=2r,b+c=2s$, with $r\ge1,s\ge2$
and $r+s=k$. Equations \eqref{eq:vander} and \eqref{eq:product} give
\[
 \N(\Th(a,b,c))\le b_rb_s\le2b_{k-1}.
\]
The legal mixed-parity triple $(1,2k-3,2)$ has count
$(k+1)b_{k-1}>2b_{k-1}$ for $k\ge3$, by \eqref{eq:short2}.

For even vertex count, all common-parity lengths are odd. Write
$c=2u+1,a+b=2v$, so $u\ge0,v\ge1,u+v=k$. The peak of the odd row
is $(2u+1)b_u/(u+1)$, and hence
\[
 \N(\Th(a,b,c))\le\frac{2u+1}{u+1}b_ub_v\le b_k.
\]
The legal mixed-parity triple $(2,2k-2,1)$ has count $kb_k>b_k$ for
$k\ge2$. At $n=4$ the all-odd class is empty, but the comparison triple
is still legal. Common-parity triples are therefore strictly dominated.
Combined with strict compression and balancing, this proves
Theorem \ref{thm:theta}.

\section{Exact maxima, ratios, and the global reduction}
Let $H(n)$ be the maximum within the simple theta class. Substitution
in \eqref{eq:short1}--\eqref{eq:short2} yields
\begin{align*}
 H(2k-1)&=\frac{k^2-2k+6}{4}\binom{k}{k/2}^{\!2}
 &&(k\text{ even}),\\
 H(2k-1)&=\frac{k(k^2-2k+5)}{k+1}
                  \binom{k-1}{(k-1)/2}^{\!2}
 &&(k\text{ odd}),\\
 H(2k)&=\frac{k^2+2}{2}\binom{k}{k/2}^{\!2}
 &&(k\text{ even}),\\
 H(2k)&=\frac{2k(k^2+1)}{k+1}
                  \binom{k-1}{(k-1)/2}^{\!2}
 &&(k\text{ odd}).
\end{align*}
For $n\ge5$, their ratios to \eqref{eq:M} are
\begin{equation}\label{eq:ratios}
\begin{array}{c|cc}
 &k\text{ even}&k\text{ odd}\\\hline
H(2k-1)/M(2k-1)&(k^2-2k+6)/(k(k+1))&(k^2-2k+5)/(k(k+1))\\[2pt]
H(2k)/M(2k)&(k^2+2)/(k(k+1))&(k^2+1)/(k(k+1)).
\end{array}
\end{equation}
All are strictly less than one in this range and tend to one along their
parity subsequences. At $n=4$, the formula gives $H(4)=12$ without any
use of $C_1$. The even-order identities also correct the intermediate
identities in \cite[Proposition 30]{BB}; the valid inequality of that
proposition was already known and is not claimed as new.

For completeness we give the block reduction. Facet-labeling shows that
$\N(C_s)=\binom{s}{s/2}$ for even $s$, and $h(s)$ for odd $s$.
It also gives multiplicativity over one-vertex sums and the factor 2
for a bridge, as in \cite[Proposition 9 and Corollary 19]{BB}.
The exact odd-cycle ratio is
\begin{equation}\label{eq:h}
 \frac{h(s+2)}{h(s)}=4\frac{s+2}{s+1}>4.
\end{equation}
This ratio decreases with $s$. Thus at fixed sum the product of two
odd-cycle counts is maximized when the lengths are balanced: transferring
two from a length $p\ge q+4$ to $q$ changes the product by
$(p-1)(q+2)/(p(q+1))>1$. This also shows $M(n+1)\ge2M(n)$: at odd
order attach a bridge, and at even order enlarge an odd cycle by two
and remove the sole bridge, gaining a factor greater than two.

Prune leaves from a graph with $n+1$ edges. Its connected 2-core with
$v$ vertices has $v+1$ edges and
\[
 \sum_{w\in V}(\deg(w)-2)=2.
\]
It has either one degree-four vertex or two degree-three vertices.
Following degree-two paths gives two cycles sharing a vertex, two
disjoint cycles joined by a path, or a theta. If its core is a theta
with $v\ge5$, its count with the $n-v$ pruned edges is
\[
 2^{n-v}\N(\Th)<2^{n-v}M(v)\le M(n).
\]
For the $v=4$ core, the count is $12\cdot2^{n-4}$, strictly below
$M(5)2^{n-5}=36\cdot2^{n-5}\le M(n)$ when $n\ge5$.

In the other cases, let the two cycle lengths be $p,q\ge3$ and let
$t$ be the number of bridges. Then
\[
 n=p+q+t-1,\qquad \N(G)=2^t\N(C_p)\N(C_q).
\]
An even cycle of length $p\ge4$ can be shortened by one and replaced
by an added bridge; its contribution improves by
$2h(p-1)/\binom p{p/2}=p/2>1$. A maximizer therefore has two odd cycles.
If $t\ge2$, remove two bridges and enlarge an odd cycle by two; by
\eqref{eq:h} the count strictly increases. Thus $t=0$ at odd order
and $t=1$ at even order. Strict balancing now gives precisely
\eqref{eq:M}. These block data are realized by simple graphs, so the
bound is attained. This proves Theorem \ref{thm:bicyclic}.

\section{Attribution, reproducibility, and review scope}
The facet-labeling theorem, cycle count, block multiplicativity,
Vandermonde identity and central-binomial ratios are existing mathematics.
Braun--Bruegge already proved the common-parity and edge-disjoint-cycle
bounds. The proposed contribution is the strict mixed-parity compression,
its full length-one boundary treatment, and the resulting exact unique
theta maximizers. The general fixed-edge result is a consequence of this
step and the stated block reduction, not a second unrelated breakthrough.

The join operation in Mori--Mori--Ohsugi \cite{MMO} adds all cross edges,
rather than identifying a single vertex. Their global join-graph theorem
does not directly give the fixed-edge theta result. The dated source
report compares that preprint and records the inaccessible
publisher-full-text boundary. No equivalent complete result was found
within the inspected material; this is not certification of historical
priority or absence of unpublished work.

The repository provides original graph-label enumeration on 19 graphs,
a separate original-graph maximizer check at orders four through eight,
and exact scalar regression on all 1,942 legal triples at orders four
through forty. The regression tests the source counting formula,
compression strictness, balancing, boundary cases, all four maxima and
their ratios. Neither finite range is a premise of a universal theorem.
The independently reviewed frozen proof sources are \texttt{PROOF.md}
and \texttt{MAXIMIZERS.md} under the workstream's \texttt{theta/}
directory. Their original source hashes and review receipts are retained.

This consolidated manuscript is an AI-assisted research draft derived
from those sources. Their separate AI review receipts do not constitute
human peer review, and this reorganization does not itself have a new
independent review receipt. No universal Lean graph/polytope theorem,
journal acceptance, DOI registration or final human authorship metadata
is represented by this file. Reproduction sources and the precise
verification map are available at
\href{https://github.com/mio-qwq/math/tree/partner/dist-STRUCT2-20261012/workstreams/STRUCT2-20261012}{the STRUCT2 workstream in the project repository}.

\begin{thebibliography}{9}
\bibitem{BB} B. Braun and K. Bruegge,
\emph{Facets of symmetric edge polytopes for graphs with few edges},
Journal of Integer Sequences \textbf{26} (2023), Article 23.7.2.
\url{https://cs.uwaterloo.ca/journals/JIS/VOL26/Braun/braun6.pdf}.
Preprint: arXiv:2201.13303v4. Conjectures 32 and 33 are journal numbering.
\bibitem{HJM} A. Higashitani, K. Jochemko and M. Micha\l ek,
\emph{Arithmetic aspects of symmetric edge polytopes},
Mathematika \textbf{65} (2019), 763--784.
\bibitem{MMO} A. Mori, K. Mori and H. Ohsugi,
\emph{Number of facets of symmetric edge polytopes arising from join graphs},
European Journal of Combinatorics \textbf{127} (2025), Article 104165.
\url{https://doi.org/10.1016/j.ejc.2025.104165}.
Inspected preprint: arXiv:2312.11287v3.
\end{thebibliography}
\end{document}
